% Chapter 1
\chapter{مقدمه}  
  
  
  \section{بیان مسئله و ضرورت پژوهش}
  
  در صنایع بزرگ، عدم اطلاع دقیق از وضعیت سلامت و عمر مفید تجهیزات، ماشین‌آلات و دستگاه‌های صنعتی، یکی از عوامل اصلی بروز خرابی‌های ناگهانی و خسارت‌های سنگین اقتصادی محسوب می‌شود. خرابی‌های غیرمنتظره در سیستم‌های حیاتی، نه تنها عملیات تولید را مختل می‌سازد، بلکه موجب افزایش هزینه‌های نگهداری، کاهش بازدهی و تهدید ایمنی کارکنان و تجهیزات نیز می‌گردد. ریشه این مشکلات را می‌توان در رویکردهای سنتی نگهداری و تعمیر دانست که بر پایه بازه‌های زمانی ثابت و بدون در نظر گرفتن شرایط واقعی تجهیزات انجام می‌شوند \cite{ismail2025systematic}.
  
دستگاه‌های صنعتی در کارخانه‌ها، نیروگاه‌ها، پالایشگاه‌ها و سامانه‌های حمل‌ونقل، معمولاً با عمر طراحی‌شده‌ای حدود ده سال تولید می‌شوند \cite{aivaliotis2021degradation}. با این حال، قرار گرفتن در معرض شرایط سخت کاری، بارگذاری‌های غیرمتعارف، تغییرات دمایی، خوردگی و فرسودگی، منجر به کاهش عمر مفید واقعی این تجهیزات می‌شود. طبق گزارش‌ها، نرخ خرابی تجهیزات صنعتی بین سال‌های ۲۰۱۰ تا ۲۰۱۸ حدود ۴۰ درصد افزایش یافته است \cite{corchado2019failure}. علاوه بر این، تنها یک روز توقف ناگهانی در فرآیند تولید می‌تواند هزینه‌ای معادل ۱۰۰٬۰۰۰ تا ۲۰۰٬۰۰۰ یورو برای سازمان به همراه داشته باشد \cite{mathur2001reasoning}.
  
در چنین شرایطی، ضرورت پژوهش در زمینه روش‌های نوین نگهداری و \gls{Monitoring} هوشمند تجهیزات بیش از پیش احساس می‌شود. یکی از فناوری‌های پیشرو در این حوزه، مفهوم «\glspl{Digital Twins}» است که با ایجاد نسخه‌ای دیجیتال از سیستم‌های فیزیکی، امکان شبیه‌سازی، پیش‌بینی و تصمیم‌گیری بهینه را فراهم می‌آورد \cite{tao2019digital}. دوقلوی دیجیتال با بهره‌گیری از داده‌های بی‌درنگ حسگرها و سامانه‌های کنترل، نمایشی پویا از وضعیت واقعی موجودیت فیزیکی ارائه می‌کند و با استفاده از مدل‌های تحلیلی و هوش مصنوعی، قادر است رفتار آینده سیستم را پیش‌بینی نماید \cite{vandinter2022predictive}.  
  
در سال‌های اخیر، با ترکیب دوقلوهای دیجیتال و یادگیری عمیق، امکان تحلیل داده‌های حجیم و پیچیده فراهم شده است. این ترکیب منجر به توسعه روش‌های پیش‌بینی خرابی، بهینه‌سازی عملکرد و تشخیص نابهنجاری‌های بی‌سابقه شده است \cite{lv2024artificial, kreuzer2024artificial}. همچنین ادغام دوقلوهای دیجیتال با \glspl{LLM} بُعد جدیدی از تعامل انسان و ماشین را ایجاد کرده است، به‌گونه‌ای که کاربران می‌توانند با زبان طبیعی با سامانه‌های صنعتی گفت‌وگو کرده و عملیات مختلفی مانند \gls{Monitoring}، تحلیل وضعیت و کنترل فرآیندها را به‌صورت هوشمند انجام دهند \cite{CHEN2025102982LLMDT}.
  
پژوهش حاضر با تمرکز بر «طراحی و پیاده‌سازی چارچوب دوقلوهای دیجیتال چندعاملی»، به دنبال توسعه یک معماری لایه‌ای و تعاملی است که در آن عامل‌ها، نقش‌های مختلفی نظیر جمع‌آوری داده، ادغام، پیش‌بینی، \gls{Monitoring} و کنترل را ایفا می‌کنند. هدف این چارچوب، ارائه راه‌حلی جامع برای مدیریت هوشمند وضعیت تجهیزات صنعتی و پیشگیری از خرابی‌های ناگهانی با بهره‌گیری از فناوری‌های نوین یادگیری عمیق و مدل‌های زبانی بزرگ است.
  


\section{اهمیت پژوهش}

پایش مستمر و پیش‌بینی خرابی تجهیزات صنعتی، یکی از ارکان اصلی افزایش قابلیت اطمینان، ایمنی و بهره‌وری در صنایع مدرن است. اهمیت این حوزه در آن است که با شناسایی زودهنگام مشکلات، از بروز خرابی‌های پرهزینه و توقف‌های ناگهانی جلوگیری می‌شود و تصمیم‌گیری در خصوص تعمیر یا جایگزینی تجهیزات بر پایه داده و تحلیل دقیق صورت می‌گیرد.

رویکردهای نگهداری صنعتی را می‌توان به سه دسته کلی تقسیم کرد \cite{ismail2025systematic}:
\begin{enumerate}
	\item \textbf{نگهداری واکنشی}: تعمیر تجهیزات پس از وقوع خرابی؛ این روش اگرچه ساده است، اما باعث توقف طولانی و هزینه‌های بالا می‌شود.
	\item \textbf{نگهداری پیشگیرانه}: انجام تعمیرات دوره‌ای در فواصل زمانی ثابت بدون توجه به وضعیت واقعی تجهیزات؛ که در بسیاری از موارد منجر به اتلاف منابع و تعویض زودهنگام قطعات می‌شود.
	\item \textbf{نگهداری پیش‌بینانه}: تحلیل مداوم داده‌های حسگرها و پیش‌بینی زمان خرابی بر اساس الگوهای واقعی عملکرد دستگاه؛ این رویکرد موجب کاهش هزینه‌ها، افزایش عمر تجهیزات و بهره‌وری بالاتر می‌شود \cite{rojek2023artificial}.
\end{enumerate}
پژوهش حاضر با تمرکز بر نگهداری پیش‌بینانه مبتنی بر دوقلوهای دیجیتال، به دنبال ایجاد بستری است که از طریق آن بتوان وضعیت سلامت تجهیزات را در هر لحظه رصد کرد و بر اساس داده‌های واقعی و الگوریتم‌های یادگیری عمیق، وقوع خرابی‌ها را پیش از رخداد شناسایی نمود. استفاده از مدل‌های زبانی بزرگ نیز امکان تعامل ساده و انسانی با این سامانه را فراهم کرده و بهره‌وری عملیات‌های صنعتی را افزایش می‌دهد.



\section{اهداف پژوهش}

اهداف این پژوهش در دو سطح کلی و اختصاصی به شرح زیر است:

\subsection*{اهداف کلی}
\begin{itemize}
	\item طراحی و توسعه یک چارچوب چندعاملی برای دوقلوهای دیجیتال به منظور پیش‌بینی خرابی و پایش بی‌درنگ تجهیزات صنعتی.
	\item ایجاد تعامل هوشمند با دوقلوی دیجیتال از طریق مدل‌های زبانی بزرگ و زبان طبیعی.
\end{itemize}

\subsection*{اهداف اختصاصی}
\begin{itemize}
	\item پیاده‌سازی دوقلوی دیجیتال برای موتورهای جت با استفاده از مجموعه داده‌های شبیه‌سازی موتورهای جت  \gls{CMAPSS} ناسا.
	\item توسعه عامل پیش‌بینی خرابی مبتنی بر شبکه عصبی \gls{LSTM}.
	\item طراحی عامل \gls{Monitoring} بی‌درنگ و عامل تعاملی مبتنی بر مدل‌های زبانی بزرگ.
	\item ارزیابی عملکرد چارچوب پیشنهادی از نظر دقت پیش‌بینی، کارایی و قابلیت تعامل.
\end{itemize}



\section{محدودیت‌های پژوهش}

\begin{itemize}
	\item
	 داده‌های استفاده‌شده در این پژوهش مربوط به مجموعه داده شبیه‌سازی‌شده \lr{CMAPSS} ناسا است که گرچه یکی از معتبرترین منابع در حوزه پیش‌بینی خرابی موتورهای جت به‌شمار می‌رود، اما داده‌های آن مصنوعی بوده و از محیط واقعی جمع‌آوری نشده است \cite{NASA2007C-MAPSS}.
	\item
	 به دلیل محدودیت دسترسی به داده‌های صنعتی واقعی، اعتبارسنجی نهایی سامانه در محیط صنعتی واقعی انجام نشده است.
	\item
	 مدل‌های زبانی بزرگ مورد استفاده برای تعامل با دوقلوی دیجیتال، نیازمند تنظیم دقیق برای دامنه صنعتی خاص هستند و در نسخه فعلی به‌صورت عمومی مورد استفاده قرار گرفته‌اند.
	\item
	 پیاده‌سازی بی‌درنگ در این پژوهش با استفاده از داده‌های شبیه‌سازی‌شده انجام شده و امکان استقرار واقعی نیازمند زیرساخت‌های خاص صنعتی است.
\end{itemize}



\section{خلاصه فصول}

ساختار این پایان‌نامه به گونه‌ای طراحی شده است که خواننده به‌صورت تدریجی از مفاهیم پایه تا پیاده‌سازی عملی و تحلیل نتایج، با مسیر پژوهش آشنا شود. بدین منظور، محتوای پایان‌نامه در پنج فصل تنظیم شده است که هر فصل نقش مشخصی در پیشبرد هدف کلی پژوهش دارد.

در \textbf{فصل نخست}، مسئله اصلی پژوهش و ضرورت پرداختن به آن در بستر صنایع مدرن تشریح می‌شود. این فصل با مرور چالش‌های موجود در زمینه نگهداری تجهیزات صنعتی و هزینه‌های ناشی از خرابی‌های ناگهانی آغاز می‌گردد و سپس اهمیت استفاده از فناوری‌های نوینی همچون دوقلوهای دیجیتال، \glspl{MAS} و مدل‌های زبانی بزرگ در بهبود فرآیند نگهداری و 
 تجهیزات بیان می‌شود. همچنین اهداف کلی و اختصاصی پژوهش، محدودیت‌های آن و ساختار کلی پایان‌نامه در این فصل معرفی می‌گردد تا مسیر پژوهش برای خواننده روشن شود.

در \textbf{فصل دوم}، مبانی نظری و ادبیات مرتبط مورد بررسی قرار گرفته است. ابتدا مفاهیم بنیادین دوقلوی دیجیتال، \glspl{MAS}، هوش مصنوعی، و یادگیری عمیق به تفصیل معرفی می‌شوند. سپس نحوه استفاده از شبکه‌های عصبی بازگشتی به‌ویژه \lr{LSTM} در مدل‌سازی زمان‌محور و پیش‌بینی خرابی تجهیزات تشریح می‌گردد. در ادامه، نقش مدل‌های زبانی بزرگ در ایجاد تعامل هوشمند با دوقلوهای دیجیتال و تسهیل تصمیم‌گیری‌های مبتنی بر زبان طبیعی توضیح داده می‌شود. بخش پایانی این فصل نیز به مرور پژوهش‌های مرتبط با دوقلوهای دیجیتال، معماری‌های چندعاملی و مجموعه داده \lr{CMAPSS} اختصاص دارد تا جایگاه پژوهش حاضر در میان مطالعات پیشین مشخص گردد.

در \textbf{فصل سوم}، روش‌شناسی پژوهش و معماری پیشنهادی ارائه می‌شود. این فصل چارچوب پنج‌لایه پیشنهادی را معرفی می‌کند که شامل \gls{EmtityLayer}، \gls{DataCollectionLayer}، \gls{IntegrationLayer}، \gls{DigitalTwinAgentsLayer} و \gls{ApplicationAgentsLayer} است. در این بخش، وظایف هر لایه، نحوه ارتباط بین عامل‌ها و \gls{Mechanism} کلی تعامل میان اجزای مختلف سامانه تشریح می‌گردد. همچنین مجموعه داده \lr{CMAPSS} ناسا به عنوان داده مبنا معرفی شده و فرآیندهای آماده‌سازی و پیش‌پردازش آن به تفصیل توضیح داده می‌شود. در پایان فصل، نحوه انطباق معماری پیشنهادی با این داده و نقش مدل‌های یادگیری عمیق و زبان طبیعی در آن بیان می‌گردد.

در \textbf{فصل چهارم}، به پیاده‌سازی عملی چارچوب پیشنهادی و آزمایش‌های انجام‌شده پرداخته می‌شود. ابتدا مراحل پیاده‌سازی و ابزارهای مورد استفاده در هر بخش تشریح شده و سپس عامل‌های مختلف سامانه شامل عامل دوقلوی دیجیتال، عامل پیش‌بینی خرابی مبتنی بر شبکه \lr{LSTM}، عامل \gls{Monitoring} بی‌درنگ و عامل تعاملی مبتنی بر مدل زبانی بزرگ به‌صورت گام‌به‌گام معرفی می‌شوند. در ادامه، نحوه شبیه‌سازی جریان داده زنده و برقراری ارتباط میان سرویس‌ها توضیح داده شده و در پایان، روش استقرار و اجرای هم‌زمان عامل‌ها در قالب یک سامانه چندعاملی ارائه می‌گردد.

در \textbf{فصل پنجم}، نتایج حاصل از آزمایش‌ها و ارزیابی عملکرد چارچوب پیشنهادی مورد تحلیل قرار می‌گیرد. در این فصل، دقت پیش‌بینی مدل \lr{LSTM}، نمودارهای عملکرد، پاسخ‌های بی‌درنگ سامانه، و نمونه‌های تعاملی میان عامل‌های زبانی و دوقلوی دیجیتال بررسی می‌شوند. همچنین میزان بهبود نسبت به روش‌های موجود، هزینه‌های محاسباتی، زمان پاسخ و قابلیت اعتماد سامانه ارزیابی می‌گردد. در ادامه، محدودیت‌های پژوهش و چشم‌اندازهای توسعه آن مطرح شده و در پایان، نوآوری‌های اصلی پژوهش شامل طراحی معماری چندلایه، ادغام مدل‌های زبانی بزرگ با دوقلوهای دیجیتال و ارائه روش تحلیل سنسورهای مؤثر بر عمر مفید تجهیزات، به‌صورت جمع‌بندی ارائه می‌شود.


\noindent
در مجموع، این پژوهش با ترکیب مفاهیم دوقلوهای دیجیتال، \glspl{MAS}، یادگیری عمیق و مدل‌های زبانی بزرگ، تلاش دارد چارچوبی نوآورانه برای مدیریت و پیش‌بینی وضعیت تجهیزات صنعتی ارائه دهد که نه‌تنها از بروز خرابی‌های ناگهانی جلوگیری کند، بلکه با فراهم کردن امکان تعامل هوشمند انسانی، گامی مهم در مسیر تحول دیجیتال صنایع بردارد.


%
%\section{اهداف پژوهش}
%
%
%
%
%هدف اصلی: طراحی چارچوب چندعاملی مبتنی بر DT برای پیش‌بینی و پایش موتورهای جت.
%اهداف فرعی: پیاده‌سازی معماری ۵ لایه‌ای، ادغام LLM برای کنترل (مانند خاموش/روشن کردن)، ارزیابی با داده‌های CMAPSS، و پیشنهاد کاربردهای صنعتی.
%اضافه کردن سؤالات تحقیق: مثلاً "چگونه ادغام MAS و LLM دقت پیش‌بینی \gls{RUL} را افزایش می‌دهد؟"
% \cite{karimi2020soft, mohammadimohammadabadi}. 
%
%
%\section{محدودیت‌های پژوهش و مجموعه داده \lr{CMAPSS}}
%
%توضیح مصنوعی بودن داده‌ها اما اعتبار بالای آن (استفاده‌شده در بیش از ۱۰۰۰ مقاله تا ۲۰۲۵)، محدودیت‌های شبیه‌سازی (عدم داده‌های واقعی بی‌درنگ)، و پتانسیل گسترش به سیستم‌های واقعی (مانند ادغام با IoT).
%
%\section{خلاصه فصل‌ها}
%توصیف مختصر ساختار پایان‌نامه.
